\documentclass[11pt,a4paper]{article}
\usepackage[T1]{fontenc}
\usepackage[utf8]{inputenc}
\usepackage{amsmath,amssymb,amsthm}
\usepackage[margin=1in]{geometry}
\usepackage[colorlinks=true,linkcolor=blue,urlcolor=blue]{hyperref}
\theoremstyle{plain}
\newtheorem{theorem}{Theorem}
\newtheorem{lemma}[theorem]{Lemma}
\newtheorem{proposition}[theorem]{Proposition}
\newtheorem{corollary}[theorem]{Corollary}
\theoremstyle{definition}
\newtheorem{remark}[theorem]{Remark}
% A quoted conjecture can be a single hundred-term formula whose terms are thirty-digit
% numbers. TeX cannot hyphenate those, so without a little stretch every such quote runs off
% the page; the linked-a-file papers are all of that shape.
\setlength{\emergencystretch}{3em}

\title{The conjectured generating function for OEIS A208995, proved}
\author{Adrian Perez Fontelles\\ \small Independent researcher}
\date{6 September 2026}
\begin{document}
\maketitle

\begin{abstract}
OEIS A208995 counts arrays of a fixed width under a local condition, and carries a conjectured
generating function --- and nothing else. It is true. The count is a walk count in a finite
digraph on $S=17$ vertices, so its generating function is a rational function whose numerator
and denominator have degrees bounded by $S$; the conjectured expression is another rational
function of known degree; and two rational functions of bounded degree agree as soon as enough
coefficients of their expansions do. Comparing $29$ coefficients in exact arithmetic
therefore settles the conjecture, rather than testing it.
\end{abstract}

\noindent\small 2020 Mathematics Subject Classification. 05A15, 05A19, 68Q45.\normalsize

\section{The conjecture}

OEIS A208995 is ``Number of 4-bead necklaces labeled with numbers -n..n not allowing reversal, with sum zero and first differences in -n..n.''. It has offset $1$ and begins
\[
2, 7, 16, 33, 58, 95,\ \dots
\]
The entry states, as a conjecture contributed by a reader and never marked settled:
\begin{quote}
G.f.: x*(2 + x - x\^{}2 + 3*x\^{}3 - x\^{}4) / ((1 - x)\^{}4*(1 + x)).
\end{quote}
As of the ``Last modified'' line on the live entry (Jul 07 2018 19:31:24, revision 16) this is still
recorded as empirical, and nothing on the entry records it as proved. The entry carries no
conjectured recurrence, which is why this generating function is the whole of what is open
here.


\section{The count is a Burnside sum}

The entry counts $4$-bead necklaces labelled from $-n..n$ with sum zero under a further
condition on consecutive beads, the beads identified under the cyclic group of rotations, written $G$. Nothing is a
walk here: the bead count is fixed and the ALPHABET grows.

By Burnside's lemma the number of orbits is $\frac{1}{|G|}\sum_{g\in G}|\mathrm{Fix}(g)|$.
A labelling fixed by $g$ is constant on the orbits of $g$, so it is a choice of one value per
orbit. For a rotation by $d$ with $c=\gcd(4,d)$ the fixed labellings are exactly the
$c$-periodic words, the whole sum is $(4/c)$ times the sum of one period, and the condition on
the $4$-periodic extension is the condition on the length-$c$ cyclic word; for a reflection
they are the palindromes, determined by $\lceil(4+1)/2\rceil$ values, and the condition is
again local in those. Every term is therefore a count of integer points of
\[
\sum_i w_i y_i = 0, \qquad -n \le y_i \le n,
\]
$w_i$ the orbit sizes, under a condition on consecutive beads. The forbidden windows the entry
names are defined by EQUALITIES between bead values, with fixed offsets, so inclusion-exclusion
over them turns each term into a count under a set of affine identifications --- one
convolution each --- and nothing is enumerated over the alphabet.

\section{The bound}

Each term is a lattice-point count over a union of relatively open rational cones in the orbit
variables, hence a quasi-polynomial in $n$; a Burnside average of quasi-polynomials is one. The
period is read off the arrangement: every ray of every cell is cut out by as many independent
hyperplanes as there are variables, so by Cramer its primitive generator has its $n$-coordinate
dividing the determinant of their $y$-parts, and a ray leaving the box or the sum-zero plane
bounds no cell and is dropped. The multiplicity of each cyclotomic factor is bounded by the
number of distinct rays whose height it divides, and by the dimension of the cone. Taking the
maximum over $g$ of the resulting multiplicities gives a monic annihilator of order
$S=17$, which includes the pre-period $n_0=0$ the conditions with an offset need. The
annihilator was then asked for terms the model had not supplied and reproduced them.

The family this entry belongs to is read by one model, described by its author as: ``Necklaces and bracelets of k beads over -n..n with sum zero AND a cyclic condition.''


The model reproduces all $44$ terms the entry publishes, exactly, in integer arithmetic:
it is built by reading the entry's own English, and a misreading gives different counts, so
this ties the model to the entry and not merely to the arithmetic.

\section{Its generating function is rational, with bounded degree}

\begin{lemma}\label{lem:rat}
$A(x)=\sum_{n\ge1}a(n)x^n=N_1(x)/D_1(x)$ where $D_1(x)$ has degree at most $S=17$ and
$D_1(0)=1$, and $\deg N_1<S+1$.
\end{lemma}

\begin{proof}
The section above derives a MONIC linear recurrence of order $S$ that $a$ provably satisfies
from some index onwards; write its characteristic polynomial as $z^S-\sum_{j<S}\alpha_jz^j$
and put $D_1(x)=1-\sum_{j<S}\alpha_jx^{S-j}$, the same polynomial written backwards, so
$\deg D_1\le S$ and $D_1(0)=1$. Multiplying the power series $A(x)$ by $D_1(x)$ annihilates
every coefficient at which the recurrence is in force, so $N_1=AD_1$ is a polynomial, and the
indices it can be supported on are the offset together with the finitely many below the point
where the recurrence takes hold, all of them less than $S+1$. No matrix is involved: the
model here is not a walk and has no adjacency matrix, and the bound comes from the recurrence
the model itself supplies.
\end{proof}

\begin{theorem}
The conjectured generating function is $A(x)$.
\end{theorem}

\begin{proof}
Write the conjecture as $N_2/D_2$ with $\deg N_2=5$ and $\deg D_2=5$; its expansion as a
power series exists because $D_2(0)\neq0$. Put
\[
R(x)\;=\;N_1(x)D_2(x)-N_2(x)D_1(x),
\]
a polynomial of degree at most $22$ by Lemma~\ref{lem:rat}. The difference
$A-N_2/D_2$ equals $R/(D_1D_2)$, and $D_1(0)D_2(0)\neq0$, so if the two power series agree in
every coefficient up to $x^{22}$ then $R$ has a zero of order greater than its own
degree at the origin, which forces $R\equiv0$ and hence $A=N_2/D_2$ identically.

The coefficients of $A$ were produced from the model by repeated matrix--vector products in
exact integer arithmetic, those of $N_2/D_2$ by exact rational long division, and $29$
of them --- more than the $22+1$ the argument requires --- agree.
\end{proof}

\section{A recurrence the entry does not state}

The denominator of a rational generating function is a linear recurrence written backwards, so
the theorem above yields one for free. With $D_2(x)=1-\sum_{i\ge1}c_ix^i$ after normalising
$D_2(0)=1$, comparing coefficients of $x^n$ in $A(x)D_2(x)=N_2(x)$ gives
\[
a(n)\;=\;\sum_{i\ge1}c_i\,a(n-i)\qquad\text{for every }n>5,
\]
a constant-coefficient linear recurrence of order $5$, valid past the degree of the
numerator. Its coefficients are the integers $c_1,\dots,c_{5}$ read off $D_2$, the first few being $3,\ -2,\ -2,\ 3,\ -1$. The entry records no recurrence, so this is not a restatement of anything
already there.

\section{Verification}

Three checks, each independent of the algebra above.

First, the model reproduces every one of the $44$ published terms, with the index shift
read off the entry's offset rather than fitted.

Second, the arrays themselves were enumerated for the smallest shapes the entry counts,
directly from its English and with no automaton, and the counts agree. Writing that check is
where the risk actually sits: an earlier version of it dropped a clause from the entry's name
in three cases and disagreed with the model, and in all three the check was wrong and the model
was right.

Third, the coefficient comparison was carried out over the whole range required by the degree
bound rather than on a prefix, in exact arithmetic throughout.

\begin{thebibliography}{9}
\bibitem{oeis} The OEIS Foundation, \emph{The On-Line Encyclopedia of Integer
Sequences}, \url{https://oeis.org}, 2026. Sequence A208995.
\bibitem{stanley} R.~P.~Stanley, \emph{Enumerative Combinatorics}, Volume 1, 2nd ed.,
Cambridge University Press, 2011. (Transfer-matrix method, Section 4.7.)
\bibitem{flajolet} P.~Flajolet and R.~Sedgewick, \emph{Analytic Combinatorics}, Cambridge
University Press, 2009.
\bibitem{horn} R.~A.~Horn and C.~R.~Johnson, \emph{Matrix Analysis}, 2nd ed., Cambridge
University Press, 2013.
\end{thebibliography}

\end{document}
